\section{Predictive entropy search}
\label{sec:method}

We propose to follow the information-theoretic method for active
data collection described in \cite{MacKay:1992}. We are interested in maximizing
information about the location $\xopt$ of the global maximum, whose posterior distribution
is $p(\xopt|\D_n)$. Our current information about $\xopt$ can
be measured in terms of the negative differential entropy of $p(\xopt|\D_n)$.
Therefore, our strategy is to select $\x_{n+1}$ which maximizes the expected
reduction in this quantity. The corresponding acquisition function is
\begin{equation}
    \textstyle
    \alpha_n(\x)
    = 
    \H[p(\xopt |\D_n)] -
    \E_{p(y|\D_n,\x)} [
        \H[p(\xopt|\D_n\cup\{(\x,y)\})]
    ]
    \,,
    \label{eq:originalObjective}
\end{equation}
where $\H[p(\x)] = -\int p(\x) \log p(\x) d\x$ represents the differential
entropy of its argument and the expectation above is taken with respect to
the posterior predictive distribution of $y$ given $\x$. 
The exact evaluation of (\ref{eq:originalObjective}) is infeasible in practice. The
main difficulties are i) $p(\xopt|\D_n\cup\{(\x,y)\})$ must be computed for
many different values of $\x$ and $y$ during the optimization of (\ref{eq:originalObjective}) and ii) the
entropy computations themselves are not analytical.  In practice, a direct
evaluation of (\ref{eq:originalObjective}) is only possible after performing
many approximations \cite{HennigSchuler:2012}.  To avoid this, we follow
the approach described in \cite{Houlsby:2012} by noting that
(\ref{eq:originalObjective}) can be equivalently written as the mutual
information between $\xopt$ and $y$ given $\D_n$. Since the mutual
information is a symmetric function, $\alpha_n(\x)$ can be rewritten as
\begin{equation}
    \textstyle
    \alpha_n(\x)
    = 
    \H[p(y|\D_n,\x)] -
    \E_{p(\xopt|\D_n)} [
        \H[p(y|\D_n,\x,\xopt)]
    ]
    \,,
    \label{eq:newObjective}
\end{equation}
where $p(y|\D_n,\x,\xopt)$ is the posterior predictive distribution for $y$
given the observed data $\D_n$ and the location of the global maximizer of $f$.
Intuitively, conditioning on the location $\xopt$ pushes the posterior predictions up in
locations around $\xopt$ and down in regions away from $\xopt$.
Note that, unlike the previous formulation, this objective is based on the
entropies of predictive distributions, which are analytic or can be easily
approximated, rather than on the entropies of distributions on $\xopt$ whose
approximation is more challenging.

The first term in (\ref{eq:newObjective}) can be computed
analytically using the posterior marginals for $f(\x)$ in (\ref{eq:gpvar}), that is,
$\H[p(y|\D_n,\x)]=0.5\log[2\pi e\,(v_n(\x)+\sigma^2)]$,
where we add $\sigma^2$ to $v_n(\x)$ because 
$y$ is obtained by adding Gaussian noise with variance $\sigma^2$ to $f(\mathbf{x})$.
The second term, on the other hand, must be approximated.
We first approximate the expectation in (\ref{eq:newObjective}) 
by averaging over samples $\xopt^{(i)}$ drawn approximately from $p(\xopt|\D_n)$.
For each of these samples, we then approximate the corresponding entropy function $\H[p(y|\D_n,\x,\xopt^{(i)})]$
using expectation propagation \cite{Minka:2001}. The code for all these operations is
publicly available at \url{http://tobediscloseduponacceptance.org}.
